\section{Effective constructions, bit resources, and scope}
We separate mathematical existence, output height, decision cost, and implementations. Let $L$ denote total binary input length, including $Q$, the rational matrix or affine equations, and the ambient dimension. All constants indexed by $n$ below are effective. No polynomial dependence on variable dimension is asserted.

\subsection{Rational support computation}
Positive semidefiniteness can be expressed by $2^n-1$ principal-minor inequalities, of degree at most $n$, in $d=n(n+1)/2$ matrix coordinates. Redundant affine equations can be removed. Safey El Din--Zhi's rational-point algorithm for convex semialgebraic sets \cite[Theorem 1.1]{SEDZ} has bit cost
\[
 \sigma^{O(1)}(sD)^{O(k^3)}
\]
for $s$ predicates of degree at most $D$ in $k$ variables and coefficient height $\sigma$, with output height $\sigma D^{O(k^3)}$. Strict inequalities are included in its quantifier-free input model. In our queries $k\leq d$, $D\leq n$, and $s\leq2^n+O(d)$. Rational projections and averaging use only matrices of fixed order and increase heights by a dimension-dependent factor. Thus at most $2n+3$ queries, together with the linear algebra between them, cost $\operatorname{poly}_n(L)$ bit operations and have output coordinate height $O_n(L)$. This is a known primitive; it is not a new general rational-point algorithm.

\subsection{The isotropy resource is not a free oracle}
For a rational nondegenerate quotient of order $k$, clear denominators to obtain an integral form whose coefficient magnitudes are at most $M=2^{O_n(L)}$. Cassels' classical small-zero theorem, in the explicit form recorded in \cite[Theorem 1.1]{Fukshansky}, supplies
\[
 0<\|v\|_\infty\leq
 \left(3\sum_{i,j}|f_{ij}|\right)^{(k-1)/2}
 \leq(3k^2M)^{(k-1)/2}
\]
whenever the integral form $v^TFv$, $F=(f_{ij})$, has a rational zero. See also \cite{CasselsSmall,CasselsBook,Dietmann}. Enumerating the finite box up to the ceiling of this bound is a total decision-and-construction procedure: every test is an exact integer evaluation. Its conservative fixed-dimensional bit cost is $2^{O_n(L+1)}$, and a successful seed has coordinate height $O_n(L+1)$. Classical local-symbol algorithms can be used instead, but their arithmetic subroutines, including any factoring assumptions, must be charged. We do not claim an unqualified deterministic polynomial-time isotropy algorithm.

The classification therefore has at most one isotropy decision in addition to the rational support queries. Once a positive answer is accompanied by a rational isotropic seed of height $h$, the remaining weighted-factor construction described next is polynomial in $L+h$ for fixed $n$. Since such a seed of height $O_n(L)$ exists, every positive instance has a polynomial-height certificate, independently of the chosen isotropy algorithm's running time.

\subsection{Constructing the exact factors}
For clarity we expand the reserve argument into effective finite steps.

In the isotropic boundary case, choose the rational hyperbolic basis of Lemma~\ref{lem:null}. In each of its charts with $a=1$ or $b=1$, the vector coordinates are polynomials of degree at most two in $z$, and the outer-product coordinates have degree at most four. Evaluate on the tensor set $\{0,1,2,3,4\}^{r-2}$. Multivariate interpolation shows that these finitely many rational null atoms span the same matrix space as the complete charts. The remaining null atoms are their limits up to scaling, so their squares lie in that span as well. Select a rational atom basis $G_i$ by elimination and put $B=\sum_iG_i$. All basis and coefficient heights are $O_n(L+h)$. Since their span contains the positive-definite-on-$V$ target, $B$ is also positive definite on $V$.

Select a positive rational $\tau$ of height $O_n(L+h)$ such that $A-\tau B\succeq0$ on $V$. This is possible from the determinant and norm bounds for rational positive-definite matrices of fixed order; alternatively successively test dyadic $\tau$ using rational principal minors. The trace condition is preserved. Lemma~\ref{lem:balance} gives a real null factor of the residual. Standard real-algebraic sampling of the equations $CC^T=A-\tau B$, $q(C_j)=0$ uses a number of variables and degree bounded in terms of $n$ only, and provides algebraic columns of bounded degree and height $O_n(L+h)$ \cite{SEDZ}. Approximate them within the rational hyperbolic charts. A nonzero algebraic chart coordinate has a root-separation lower bound $2^{-O_n(L+h)}$; if both hyperbolic coordinates vanish, first perturb by a small rational multiple of $v$ as in Lemma~\ref{lem:null}. Thus the rational parameters and atom approximations needed for accuracy $2^{-p}$ have height $O_n(L+h+p)$.

The inverse matrix of the selected rational atom basis has height $O_n(L+h)$ by Cramer's rule. Taking $p=O_n(L+h)$ therefore makes the residual basis-coefficient error smaller than $\tau/2$. All corrected coefficients are positive, the trace equality remains exact, and the reconstructed matrix is exactly $A$. Rational dependence elimination reduces to at most $r(r+1)/2-1$ weighted atoms.

In the positive-trace case, find a rational vector with $q>0$ by rational diagonalization. A rational box around it remains in $q>0$, with radius of height $O_n(L)$ obtainable by direct quadratic estimates. Evaluating rank-one squares on a degree-two interpolation grid in this box gives a rational atom basis of $\Sym(V)$. Choose $\tau$ so that $A-\tau B$ stays positive definite on $V$ and has positive trace. Sample a real factor with strictly positive column values; fixed-dimensional real-algebraic height bounds give sufficient positive clearance. Rational rounding and the same reserve repair produce at most $r(r+1)/2$ weighted atoms of height $O_n(L)$. If $q|_V=0$, ordinary rational diagonalization suffices. These constructions also apply to the rational reserve matrix returned by the affine-carrier algorithm.

A positive rational coefficient $a/b$ can be converted to at most four rational squares with coordinate height $O(\operatorname{bits}(a)+\operatorname{bits}(b))$. This is an existence bound. For a deterministic polynomial-time implementation avoiding a four-square search, expand the integer $ab$ in binary: an even power of two is a square, and an odd power is a sum of two equal squares. Dividing by $b$ gives $O(\operatorname{bits}(ab))$ rational squares. Therefore deterministic conversion may output $O_n(L+h)$ columns rather than the four-times-weighted-width existence bound.

\subsection{Finite precision and observations}
A rational target in the terminal slice can be approximated within $2^{-p}$ with its affine measurements unchanged by moving a rational distance toward $X_*$ and applying the preceding construction. Fixed-dimensional determinant estimates control the new relative-interior margin; output heights are $O_n(L+h+p)$ and post-seed construction costs are polynomial in those quantities. This is not an exact membership claim for an unperturbed target outside the rational-factor set.

For a full matrix outside $C_q$, the cone is closed, so real separation applies. When the real fiber is feasible but rationally impossible, no real separating form can certify that arithmetic impossibility. For a nonclosed observation image, even real infeasibility can lack strict separation. The carrier theorem and the local quadratic-form obstruction give the correct alternatives rather than conflating these cases.

\subsection{Verification and the relation to prior local-pivot obstructions}
The accompanying standard-library scripts verify exact low-rank obstruction matrices, their range restrictions and finite modular obstructions; rational null and strict factors; an algebraic totally-real factor identity; signed Cayley identities; and homogeneous circuit blocks. A separate verifier imports no generator code and rejects deliberately corrupted data. These are checks of instances and interfaces. Neither a production general rational-point solver nor a general isotropy implementation is claimed.

The earlier Manuscript II branch proves exponential prefixes for every strictly descending extreme-ray adjacent rule on robust nonmembers of a fixed nonsimplicial column cone. The present theorem does not transfer that lower bound to every Lorentz cone, every initialization, or every algorithm. It instead answers a different generality question: the arithmetic content of a real boundary factorization depends on its quadratic zero domain. A short description and good interior conditioning do not alone control its rational feasibility, and rational certificate feasibility does not alone control a chosen local adjacency graph. The product reduction gives an independent reason why a proposed universal scalar condition for all matrix-cone problems would need qualifications.
